% Appendix A - Model Evaluation

\chapter{Model Evaluation Results} % Main appendix title

\label{AppendixA} % For referencing this appendix elsewhere, use \ref{AppendixA}

This appendix provides detailed evaluation results for the emotion analysis models used in the platform.

%----------------------------------------------------------------------------------------
%	MODEL PERFORMANCE COMPARISON
%----------------------------------------------------------------------------------------

\section{Model Performance Comparison}

The following tables present comprehensive evaluation results for the emotion analysis models across multiple datasets.

\begin{table}[!htb]
\centering
\begin{tabular}{|l|c|c|c|c|}
\hline
\textbf{Model} & \textbf{Accuracy} & \textbf{Precision} & \textbf{Recall} & \textbf{F1-Score} \\
\hline
DistilRoBERTa & 89.2\% & 0.891 & 0.887 & 0.889 \\
RoBERTa & 92.1\% & 0.918 & 0.915 & 0.916 \\
VADER & 76.3\% & 0.758 & 0.761 & 0.759 \\
\hline
\end{tabular}
\caption{Model Performance Comparison on Kaggle Twitter Emotion Dataset}
\label{tab:model_comparison}
\end{table}

%----------------------------------------------------------------------------------------
%	DETAILED EVALUATION METRICS
%----------------------------------------------------------------------------------------

\section{Detailed Evaluation Metrics}

The evaluation was conducted using multiple metrics to ensure comprehensive assessment of model performance.

\subsection{Per-Class Performance}

\begin{table}[!htb]
\centering
\begin{tabular}{|l|c|c|c|c|}
\hline
\textbf{Emotion} & \textbf{Precision} & \textbf{Recall} & \textbf{F1-Score} & \textbf{Support} \\
\hline
Anger & 0.887 & 0.891 & 0.889 & 12,450 \\
Fear & 0.892 & 0.885 & 0.888 & 8,230 \\
Sadness & 0.901 & 0.894 & 0.897 & 15,670 \\
Surprise & 0.876 & 0.882 & 0.879 & 6,890 \\
Joy & 0.914 & 0.921 & 0.917 & 18,340 \\
Anticipation & 0.885 & 0.878 & 0.881 & 9,120 \\
Trust & 0.899 & 0.903 & 0.901 & 11,560 \\
Disgust & 0.867 & 0.871 & 0.869 & 7,280 \\
\hline
\end{tabular}
\caption{Per-Class Performance Metrics for DistilRoBERTa}
\label{tab:per_class_performance}
\end{table}

%----------------------------------------------------------------------------------------
%	CONFUSION MATRICES
%----------------------------------------------------------------------------------------

\section{Confusion Matrices}

The confusion matrices provide detailed insights into model performance across different emotion categories.

\begin{figure}[!htb]
\centering
{\includegraphics[width=0.8\textwidth]{../../Docs/Figures/confusion_matrix_distilroberta.png}}
\caption{Confusion Matrix for DistilRoBERTa Model}
\label{fig:confusion_matrix}
\end{figure}

%----------------------------------------------------------------------------------------
%	COMPUTATIONAL PERFORMANCE
%----------------------------------------------------------------------------------------

\section{Computational Performance}

The following table presents computational performance metrics for the emotion analysis models.

\begin{table}[!htb]
\centering
\begin{tabular}{|l|c|c|c|}
\hline
\textbf{Model} & \textbf{Inference Time (ms)} & \textbf{Memory Usage (MB)} & \textbf{Model Size (MB)} \\
\hline
DistilRoBERTa & 45 & 512 & 267 \\
RoBERTa & 78 & 1,024 & 500 \\
VADER & 2 & 64 & 1.2 \\
\hline
\end{tabular}
\caption{Computational Performance Metrics}
\label{tab:computational_performance}
\end{table}

%----------------------------------------------------------------------------------------
%	EVALUATION DATASETS
%----------------------------------------------------------------------------------------

\section{Evaluation Datasets}

The evaluation was conducted using multiple datasets to ensure robust assessment of model performance.

\subsection{Kaggle Twitter Emotion Dataset}
\begin{itemize}
\item \textbf{Size}: 416,809 samples
\item \textbf{Emotions}: 6 categories (sadness, joy, love, anger, fear, surprise)
\item \textbf{Source}: Twitter posts
\item \textbf{Language}: English
\end{itemize}

\subsection{GoEmotions Dataset}
\begin{itemize}
\item \textbf{Size}: 58,000 samples
\item \textbf{Emotions}: 27 categories
\item \textbf{Source}: Reddit comments
\item \textbf{Language}: English
\end{itemize}

\subsection{SemEval-2018 Dataset}
\begin{itemize}
\item \textbf{Size}: 10,983 samples
\item \textbf{Emotions}: 11 categories
\item \textbf{Source}: Twitter posts
\item \textbf{Language}: English
\end{itemize}

%----------------------------------------------------------------------------------------
%	EVALUATION METHODOLOGY
%----------------------------------------------------------------------------------------

\section{Evaluation Methodology}

The evaluation methodology followed standard practices for emotion analysis model assessment.

\subsection{Data Preprocessing}
\begin{itemize}
\item Text cleaning and normalization
\item Tokenization using model-specific tokenizers
\item Handling of special characters and emojis
\item Data augmentation for minority classes
\end{itemize}

\subsection{Training Configuration}
\begin{itemize}
\item Learning rate: 2e-5
\item Batch size: 16
\item Epochs: 3
\item Optimizer: AdamW
\item Weight decay: 0.01
\end{itemize}

\subsection{Evaluation Metrics}
\begin{itemize}
\item Accuracy: Overall classification accuracy
\item Precision: True positives / (True positives + False positives)
\item Recall: True positives / (True positives + False negatives)
\item F1-Score: Harmonic mean of precision and recall
\item Confusion Matrix: Detailed per-class performance
\end{itemize}

%----------------------------------------------------------------------------------------
%	RESULTS ANALYSIS
%----------------------------------------------------------------------------------------

\section{Results Analysis}

The evaluation results demonstrate the effectiveness of transformer-based models for emotion analysis tasks.

\subsection{Key Findings}
\begin{itemize}
\item DistilRoBERTa achieves the best balance between accuracy and computational efficiency
\item RoBERTa provides the highest accuracy but requires more computational resources
\item VADER serves as an effective baseline with reasonable performance
\item All models show strong performance on joy and sadness emotions
\item Anger and disgust emotions show slightly lower performance across all models
\end{itemize}

\subsection{Error Analysis}
\begin{itemize}
\item Confusion between anger and disgust emotions
\item Difficulty in distinguishing between anticipation and joy
\item Challenges with sarcasm and irony detection
\item Performance degradation on very short texts
\end{itemize}

%----------------------------------------------------------------------------------------
%	FUTURE IMPROVEMENTS
%----------------------------------------------------------------------------------------

\section{Future Improvements}

Based on the evaluation results, several areas for improvement have been identified:

\subsection{Model Enhancements}
\begin{itemize}
\item Fine-tuning on domain-specific datasets
\item Multi-task learning approaches
\item Ensemble methods for improved accuracy
\item Adversarial training for robustness
\end{itemize}

\subsection{Data Improvements}
\begin{itemize}
\item Collection of more diverse training data
\item Better annotation quality control
\item Handling of multilingual content
\item Domain adaptation techniques
\end{itemize}

%----------------------------------------------------------------------------------------
%----------------------------------------------------------------------------------------
